\documentclass[12pt,letterpaper]{report}

\usepackage[left=3cm,right=3cm,top=3cm,bottom=3cm]{geometry}
\usepackage{times}
\usepackage[T1]{fontenc}
\usepackage[utf8]{inputenc}
\usepackage{graphicx}
\usepackage{float}
\usepackage{enumitem}
\usepackage{hyperref}
\usepackage{titlesec}
\usepackage{setspace}
\usepackage{listings}
\usepackage{xcolor}
\usepackage{booktabs}
\usepackage{longtable}

\hypersetup{colorlinks=true, urlcolor=blue, linkcolor=black}
\onehalfspacing
\setlength{\parindent}{0pt}
\setlength{\parskip}{6pt}

\titleformat{\section}{\normalfont\bfseries}{\Alph{section}.}{0.5em}{}
\titleformat{\subsection}{\normalfont\bfseries}{\arabic{subsection})}{0.5em}{}
\titleformat{\subsubsection}{\normalfont\itshape}{}{0em}{}

\lstset{
    basicstyle=\ttfamily\small,
    breaklines=true,
    frame=single,
    backgroundcolor=\color{gray!10},
    columns=flushleft
}

\newcommand{\shotwidth}{0.75\textwidth}
\newcommand{\shot}[3]{%
  \begin{figure}[H]
    \centering
    \IfFileExists{images/#1}{%
      \includegraphics[width=\shotwidth]{images/#1}%
    }{%
      \fbox{%
        \begin{minipage}[c][5cm][c]{\shotwidth}
          \centering
          \textit{[Placeholder Screenshot -- ganti dengan hasil pengerjaan Anda]}\\[6pt]
          \url{images/#1}
        \end{minipage}%
      }%
    }
    \caption{#2}
    \label{fig:#3}
  \end{figure}
}

% Tandai bagian yang HARUS kamu isi sendiri dengan hasil nyata
\newcommand{\isi}[1]{\textbf{\textcolor{red}{[ISI: #1]}}}

\begin{document}

\begin{titlepage}
\centering
\vspace*{0.3cm}
{\large \textbf{LAPORAN UJIAN TENGAH SEMESTER}}\\[4pt]
{\large \textbf{KONSTRUKSI \& EVOLUSI PERANGKAT LUNAK}}\\[4pt]
{\large \textbf{LARAVEL, VUE, CI/CD GITHUB ACTIONS, DOCKER, DAN DOCKER COMPOSE}}

\vspace{1.2cm}
\includegraphics[width=3.5cm]{images/lambang ugm.png}
\vspace{1.2cm}

\begin{minipage}{0.75\textwidth}
\centering
\begin{tabular}{@{}p{3.2cm}p{0.4cm}l@{}}
\multicolumn{3}{l}{Disusun oleh:} \\[6pt]
Nama            & : & Mohammad Avirza Radyatanza \\[8pt]
NIM             & : & 24/534267/SV/24022 \\[8pt]
Kelas           & : & PL5AA \\[8pt]
Dosen Pengampu  & : & Galih Malela Damaraji, S.Pd., M.Eng. \\
\end{tabular}
\end{minipage}

\vspace{1.2cm}

{\large \textbf{PROGRAM STUDI D-IV TEKNOLOGI REKAYASA PERANGKAT LUNAK}}\\[4pt]
{\large \textbf{DEPARTEMEN TEKNIK ELEKTRO DAN INFORMATIKA}}\\[4pt]
{\large \textbf{SEKOLAH VOKASI}}\\[4pt]
{\large \textbf{UNIVERSITAS GADJAH MADA}}\\[4pt]
{\large \textbf{YOGYAKARTA}}\\[4pt]
{\large \textbf{2026}}

\end{titlepage}

\section{Informasi Umum}
\begin{itemize}[label=--]
    \item Tema sistem: \textbf{Project Tracker / Todo Management} -- Aplikasi manajemen proyek dan tugas untuk seorang \textit{Project Manager} yang membawahi banyak anggota tim. PM dapat membuat proyek, menugaskan tugas ke anggota, dan memantau kemajuan.
    \item Repositori GitHub: \url{https://github.com/MohammadAvirzaR/uts-pl-24-534267-SV-24022}
    \item Halaman GitHub Actions: \url{https://github.com/MohammadAvirzaR/uts-pl-24-534267-SV-24022/actions}
    \item Halaman Packages (GHCR): \url{https://github.com/MohammadAvirzaR/packages}
    \item Postman Collection: berkas \texttt{backend/postman/ProjectTracker.postman\_collection.json} pada repositori
\end{itemize}

% ============================================================
%  SOAL 1 -- BACKEND LARAVEL (CPMK1, 15pt)
% ============================================================
\section{Soal 1 -- Backend Laravel RESTful API (CPMK1)}

\subsection{Autentikasi Token (Login dan Logout)}
Autentikasi menggunakan \textbf{Laravel Sanctum}, paket bawaan Laravel 11 yang menyediakan sistem token sederhana untuk SPA dan API. Proses login dilakukan melalui endpoint \texttt{POST /api/login} yang memvalidasi email dan password, lalu membuat token menggunakan \texttt{createToken()}. Token dikembalikan dalam response JSON dan harus disertakan di header \texttt{Authorization: Bearer <token>} pada setiap request selanjutnya.

Logout dilakukan melalui \texttt{POST /api/logout} yang mencabut (\textit{revoke}) token yang sedang digunakan melalui \texttt{currentAccessToken()->delete()}. Seluruh endpoint selain \texttt{/api/login} dan \texttt{/api/register} dilindungi oleh middleware \texttt{auth:sanctum}, sehingga mengembalikan \textbf{401 Unauthorized} apabila request tidak menyertakan token valid.

\begin{lstlisting}[language=php,caption={Potongan kode AuthController -- login}]
public function login(Request $request)
{
    $request->validate([
        'email' => 'required|email',
        'password' => 'required',
    ]);

    $user = User::where('email', $request->email)->first();

    if (!$user || !Hash::check($request->password, $user->password)) {
        throw ValidationException::withMessages([
            'email' => ['The provided credentials are incorrect.'],
        ]);
    }

    $token = $user->createToken('auth-token')->plainTextToken;

    return response()->json([
        'message' => 'Login successful',
        'user' => $user,
        'token' => $token,
    ]);
}
\end{lstlisting}

\shot{1.png}{Postman: login berhasil dan token diterima}{s1-login}
\shot{2.png}{Postman: akses tanpa token mengembalikan 401}{s1-401}

\subsection{Resource CRUD}
Dua resource utama disediakan: \texttt{Project} dan \texttt{Task}. Berikut tabel endpoint lengkap:

\begin{table}[H]
\centering
\begin{tabular}{lllll}
\toprule
\textbf{Method} & \textbf{Endpoint} & \textbf{Fungsi} & \textbf{Status} \\
\midrule
POST   & /api/register         & Registrasi user baru      & 201 \\
POST   & /api/login            & Login, dapatkan token      & 200 \\
POST   & /api/logout           & Logout, cabut token        & 200 \\
GET    & /api/user             & Profil user saat ini       & 200 \\
\midrule
GET    & /api/projects         & Tampil semua project       & 200 \\
POST   & /api/projects         & Buat project baru          & 201 \\
GET    & /api/projects/\{id\}  & Detail satu project        & 200 \\
PUT    & /api/projects/\{id\}  & Ubah project               & 200 \\
DELETE & /api/projects/\{id\}  & Hapus project              & 200 \\
\midrule
GET    & /api/tasks            & Tampil semua task          & 200 \\
POST   & /api/tasks            & Buat task baru             & 201 \\
GET    & /api/tasks/\{id\}     & Detail satu task           & 200 \\
PUT    & /api/tasks/\{id\}     & Ubah task                  & 200 \\
DELETE & /api/tasks/\{id\}     & Hapus task                 & 200 \\
\bottomrule
\end{tabular}
\caption{Daftar endpoint RESTful API}
\end{table}

\begin{lstlisting}[language=php,caption={Potongan kode ProjectController -- store}]
public function store(Request $request)
{
    $validated = $request->validate([
        'name' => 'required|string|max:255',
        'description' => 'nullable|string',
        'status' => 'sometimes|in:planning,active,completed,on_hold',
        'deadline' => 'nullable|date',
    ]);

    $project = $request->user()->projects()->create($validated);

    return response()->json([
        'message' => 'Project created successfully',
        'data' => $project,
    ], 201);
}
\end{lstlisting}

\shot{3.png}{Postman: GET daftar data}{s1-crud-index}
\shot{4.png}{Postman: POST data baru, status 201}{s1-crud-store}
\shot{5.png}{Postman: validasi gagal, status 422}{s1-crud-422}
\shot{6.png}{Postman: PUT/PATCH dan DELETE}{s1-crud-ud}

\subsection{Feature Test}
Dua file feature test dibuat:

\begin{enumerate}[label=\alph*.]
    \item \texttt{AuthTest.php} -- Menguji: login berhasil mengembalikan token, akses tanpa token mengembalikan 401, dan logout berhasil mencabut token.
    \item \texttt{ProjectTest.php} -- Menguji: list project (200), create project (201), validasi gagal (422), update project (200), dan delete project (200).
\end{enumerate}

\begin{lstlisting}[language=php,caption={Potongan AuthTest.php}]
public function test_user_can_login(): void
{
    $user = User::factory()->create([
        'password' => Hash::make('password123'),
    ]);

    $response = $this->postJson('/api/login', [
        'email' => $user->email,
        'password' => 'password123',
    ]);

    $response->assertStatus(200)
             ->assertJsonStructure(['token', 'user']);
}

public function test_user_cannot_access_protected_route_without_token(): void
{
    $response = $this->getJson('/api/projects');
    $response->assertStatus(401);
}
\end{lstlisting}

\shot{7.png}{Output php artisan test (lulus)}{s1-test}

\subsection{Postman Collection}
Seluruh endpoint diekspor sebagai Postman Collection dan disimpan di direktori \texttt{backend/postman/} dalam repositori.

\shot{8.png}{Berkas Postman Collection di dalam repositori}{s1-postman-repo}

% ============================================================
%  SOAL 2 -- FRONTEND VUE.JS (CPMK1, 10pt)
% ============================================================
\section{Soal 2 -- Frontend Vue.js (CPMK1)}

\subsection{Halaman Login dan Penyimpanan Token}
Token dari response login disimpan di \texttt{localStorage} melalui Pinia store (\texttt{auth.js}). Sebuah Axios interceptor dikonfigurasi untuk secara otomatis menyertakan header \texttt{Authorization: Bearer <token>} pada setiap request, sehingga komponen tidak perlu menangani token secara manual.

\begin{lstlisting}[language=python]
import axios from 'axios'

const api = axios.create({
    baseURL: import.meta.env.VITE_API_URL || 'http://localhost:8000/api',
})

api.interceptors.request.use((config) => {
    const token = localStorage.getItem('token')
    if (token) {
        config.headers.Authorization = `Bearer ${token}`
    }
    return config
})

export default api
\end{lstlisting}

\shot{9.png}{Halaman login Vue}{s2-login}

\subsection{Halaman CRUD melalui Axios}
Halaman utama menampilkan daftar project milik user. Dari halaman tersebut, user dapat melihat detail, membuat baru, mengedit, dan menghapus project. Setiap project juga memiliki halaman detail yang menampilkan task-task di dalamnya.

\shot{10.png}{Halaman tampil data}{s2-list}
\shot{11.png}{Form tambah dan ubah data}{s2-form}

\subsection{Alamat API melalui Variabel Lingkungan}
Alamat API backend \textbf{tidak ditulis mati} (hardcode) di dalam kode komponen. Alamat tersebut dikonfigurasi melalui variabel \texttt{VITE\_API\_URL} pada berkas \texttt{.env} frontend. Variabel ini dibaca melalui \texttt{import.meta.env.VITE\_API\_URL} oleh Vite pada saat build. Hal ini memudahkan pergantian alamat API tanpa menyunting kode sumber, misalnya saat berpindah dari environment development ke Docker Compose.

\begin{lstlisting}[language=bash,caption={File .env frontend}]
VITE_API_URL=/api
\end{lstlisting}

% ============================================================
%  SOAL 3 -- GIT & CI (CPMK2, 15pt)
% ============================================================
\section{Soal 3 -- Manajemen GitHub dan Continuous Integration (CPMK2)}

\subsection{Struktur Repositori dan Branch}
Repository publik dibuat dengan nama \texttt{uts-pl-24-534267-SV-24022-ver1}, berisi dua folder: \texttt{backend/} untuk Laravel dan \texttt{frontend/} untuk Vue. Tiga branch digunakan: \texttt{main} (produksi), \texttt{dev} (integrasi), dan branch \texttt{feature/*} untuk pengembangan fitur.

\shot{12.png}{Tab Code repositori: folder backend/ dan frontend/}{s3-repo}
\shot{13.png}{Daftar branch main, dev, dan branch fitur}{s3-branches}

\subsection{Pull Request}
Semua perubahan masuk ke \texttt{dev} dan \texttt{main} melalui Pull Request, tidak pernah push langsung.

\shot{14.png}{Pull Request yang di-merge ke dev dan main}{s3-pr}

\subsection{Branch Protection}
Branch protection diaktifkan pada \texttt{main} dan \texttt{dev} melalui Settings $\rightarrow$ Branches (atau Rules). Require Pull Request dan status check lulus sebelum merge.

\shot{15.png}{Pengaturan branch protection (Settings $\rightarrow$ Branches/Rules)}{s3-protection}

\subsection{Workflow CI}
Workflow \texttt{.github/workflows/ci.yml} berisi minimal dua job:
\begin{enumerate}[label=\alph*.]
    \item \textbf{test}: Menjalankan \texttt{php artisan test} untuk menguji backend.
    \item \textbf{build}: Menjalankan \texttt{npm ci} dan \texttt{npm run build} untuk memvalidasi frontend.
\end{enumerate}

Workflow berjalan pada setiap push dan pull request ke branch \texttt{main} dan \texttt{dev}.

\begin{lstlisting}[language=python]
name: CI/CD Pipeline
on:
  push:
    branches: [main, dev]
  pull_request:
    branches: [main, dev]

jobs:
  test:
    runs-on: ubuntu-latest
    steps:
      - uses: actions/checkout@v4
      - name: Setup PHP
        uses: shivammathur/setup-php@v2
        with: { php-version: '8.3' }
      - run: composer install --no-progress
        working-directory: backend
      - run: cp .env.example .env && php artisan key:generate
        working-directory: backend
      - run: php artisan test
        working-directory: backend

  build:
    runs-on: ubuntu-latest
    steps:
      - uses: actions/checkout@v4
      - uses: actions/setup-node@v4
        with: { node-version: '22' }
      - run: npm ci
        working-directory: frontend
      - run: npm run build
        working-directory: frontend
\end{lstlisting}

\shot{16.png}{Workflow CI berstatus hijau pada tab Actions}{s3-ci-hijau}

% ============================================================
%  SOAL 4 -- CD (CPMK3, 15pt)
% ============================================================
\section{Soal 4 -- Continuous Deployment (CPMK3)}

\subsection{Kondisi Job Deploy}
Job \texttt{deploy} ditambahkan pada workflow yang sama, dengan kondisi \texttt{if: github.ref == 'refs/heads/main' \&\& github.event\_name == 'push'} sehingga hanya berjalan saat ada push ke branch \texttt{main}. Job ini bergantung (\texttt{needs}) pada job \texttt{test} dan \texttt{build}, sehingga deploy hanya terjadi setelah pengujian dan build berhasil.

\subsection{GitHub Secrets}
Nilai rahasia \texttt{APP\_KEY} dan \texttt{DB\_PASSWORD} disimpan di GitHub Secrets (Settings $\rightarrow$ Secrets and variables $\rightarrow$ Actions). Nilai-nilai ini \textbf{tidak} ada di dalam kode repository.

\shot{17.png}{Daftar GitHub Secrets (nilai tidak terlihat)}{s4-secrets}

\subsection{Bukti Run}
\shot{18.png}{Run pada branch main dengan job deploy hijau}{s4-main}
\shot{19.png}{Run dari dev atau Pull Request: job deploy tidak berjalan}{s4-dev}
\shot{20.png}{Halaman Packages berisi image backend dan frontend}{s4-packages}

% ============================================================
%  SOAL 5 -- DOCKERFILE (CPMK4 & CPMK5, 20pt)
% ============================================================
\section{Soal 5 -- Dockerfile (CPMK4 dan CPMK5)}

\subsection{Dockerfile Backend}
Backend menggunakan image dasar \texttt{dunglas/frankenphp:1.4-php8.4} (versi eksplisit, bukan \texttt{latest}). File \texttt{composer.json} dan \texttt{composer.lock} disalin \textbf{terlebih dahulu} sebelum kode aplikasi, lalu \texttt{composer install} dijalankan. Dengan demikian, selama dependensi tidak berubah, layer cache Composer dapat digunakan kembali dan proses build menjadi jauh lebih cepat.

\begin{lstlisting}[language=bash,caption={backend/Dockerfile (FrankenPHP)}]
FROM dunglas/frankenphp:1.4-php8.4

RUN install-php-extensions pdo_pgsql pcntl

WORKDIR /app

# Layer cache: install dependencies first
COPY composer.json composer.lock ./
RUN composer install --no-dev --no-scripts --prefer-dist

# Then copy app code
COPY . .
RUN composer dump-autoload --optimize

RUN php artisan config:cache \
    && php artisan route:cache

EXPOSE 8000

CMD ["php", "artisan", "octane:frankenphp", "--host=0.0.0.0", "--port=8000"]
\end{lstlisting}

\subsection{Dockerfile Frontend (Multi-Stage)}
Frontend menggunakan \textit{multi-stage build} dengan dua tahap:
\begin{enumerate}[label=\arabic*.]
    \item \textbf{Tahap build} (\texttt{node:22-alpine}): menjalankan \texttt{npm ci} dan \texttt{npm run build} untuk menghasilkan file statis di \texttt{dist/}.
    \item \textbf{Tahap runtime} (\texttt{nginx:1.27-alpine}): hanya menyalin hasil build dari tahap pertama. Image akhir \textbf{tidak berisi} \texttt{node\_modules}.
\end{enumerate}

\begin{lstlisting}[language=bash,caption={frontend/Dockerfile (Multi-stage)}]
# Stage 1: Build
FROM node:22-alpine AS build
WORKDIR /app
COPY package.json package-lock.json ./
RUN npm ci
COPY . .
RUN npm run build

# Stage 2: Serve
FROM nginx:1.27-alpine AS runtime
COPY --from=build /app/dist /usr/share/nginx/html
COPY nginx.conf /etc/nginx/conf.d/default.conf
EXPOSE 80
\end{lstlisting}

\subsection{.dockerignore dan Tidak Ada Rahasia}
Berkas \texttt{.dockerignore} dibuat di masing-masing folder untuk mengecualikan:
\begin{itemize}[nosep]
    \item \texttt{backend/}: \texttt{vendor/}, \texttt{.env}, \texttt{storage/logs/}, \texttt{node\_modules/}, \texttt{tests/}
    \item \texttt{frontend/}: \texttt{node\_modules/}, \texttt{dist/}, \texttt{.env}, \texttt{.env.local}
\end{itemize}

Tidak ada nilai rahasia (password, key, token) yang ditulis di dalam Dockerfile. Seluruh konfigurasi sensitif diinjeksi melalui variabel lingkungan saat runtime via Docker Compose.

\shot{21.png}{Output docker images: ukuran image akhir dan tanpa node\_modules}{s5-images}

% ============================================================
%  SOAL 6 -- DOCKER COMPOSE (CPMK6 & CPMK7, 25pt)
% ============================================================
\section{Soal 6 -- Docker Compose (CPMK6 dan CPMK7)}

\subsection{Berkas docker-compose.yml}
Satu berkas \texttt{docker-compose.yml} di akar repository mendefinisikan tiga service:

\begin{enumerate}[label=\alph*.]
    \item \textbf{db}: PostgreSQL 16 dengan named volume \texttt{pgdata} untuk persistensi data.
    \item \textbf{backend}: Laravel + FrankenPHP, terhubung ke \texttt{db} melalui nama service Docker.
    \item \textbf{frontend}: Vue + Nginx, meneruskan request API ke \texttt{backend} melalui nama service.
\end{enumerate}

Komunikasi antar service menggunakan \textbf{nama service di jaringan Docker}, bukan alamat IP laptop.

\begin{lstlisting}[language=python]
services:
  db:
    image: postgres:16-alpine
    environment:
      POSTGRES_DB: ${DB_DATABASE}
      POSTGRES_USER: ${DB_USERNAME}
      POSTGRES_PASSWORD: ${DB_PASSWORD}
    volumes:
      - pgdata:/var/lib/postgresql/data
    healthcheck:
      test: ["CMD-SHELL", "pg_isready -U ${DB_USERNAME}"]
      interval: 5s
      timeout: 5s
      retries: 5

  backend:
    build: ./backend
    environment:
      APP_KEY: ${APP_KEY}
      DB_CONNECTION: pgsql
      DB_HOST: db
      DB_PORT: 5432
      DB_DATABASE: ${DB_DATABASE}
      DB_USERNAME: ${DB_USERNAME}
      DB_PASSWORD: ${DB_PASSWORD}
    ports:
      - "8000:8000"
    depends_on:
      db:
        condition: service_healthy

  frontend:
    build: ./frontend
    ports:
      - "5173:80"
    depends_on:
      - backend

volumes:
  pgdata:
\end{lstlisting}

\subsection{Menjalankan Sistem}
\begin{lstlisting}[language=bash]
docker compose up -d --build
\end{lstlisting}

\shot{22.png}{Output docker compose up -d --build}{s6-up}
\shot{23.png}{Output docker compose ps}{s6-ps}

\subsection{Pengujian dari Browser}
\shot{24.png}{Login dari browser pada sistem Compose}{s6-browser-login}
\shot{25.png}{CRUD dari browser; data tersimpan di basis data}{s6-browser-crud}

\subsection{Persistensi Data}
Data basis data disimpan di \textit{named volume} \texttt{pgdata}. Setelah menjalankan \texttt{docker compose down} (tanpa \texttt{-v}) kemudian \texttt{docker compose up -d}, data yang sudah dibuat sebelumnya \textbf{tetap ada} dan tidak hilang.

\shot{26.png}{Data tetap ada setelah docker compose down lalu up -d}{s6-volume}

\subsection{Kredensial dari .env}
Seluruh kredensial basis data dibaca dari berkas \texttt{.env} yang \textbf{tidak} di-commit ke repository (tercantum di \texttt{.gitignore}). Berkas \texttt{.env.example} disertakan sebagai contoh konfigurasi:

\begin{lstlisting}[language=bash,caption={.env.example}]
APP_KEY=base64:your-key-here
DB_DATABASE=project_tracker
DB_USERNAME=postgres
DB_PASSWORD=secret
\end{lstlisting}

% ============================================================
%  SOAL 7 -- BONUS FRANKENPHP (+10pt)
% ============================================================
\section{Soal 7 -- Nilai Tambahan: FrankenPHP}

Backend Laravel menggunakan \textbf{FrankenPHP} sebagai server aplikasi, menggantikan \texttt{php artisan serve} atau PHP-FPM dengan Nginx. Image dasar yang digunakan adalah \texttt{dunglas/frankenphp:1.4-php8.4} (versi eksplisit, bukan \texttt{latest}).

FrankenPHP dijalankan melalui Laravel Octane dengan perintah:
\begin{lstlisting}[language=bash]
php artisan octane:frankenphp --host=0.0.0.0 --port=8000
\end{lstlisting}

Backend tetap berjalan bersama service lain (\texttt{db} dan \texttt{frontend}) di dalam \texttt{docker-compose.yml}.

\shot{27.png}{Log container backend menunjukkan permintaan dilayani FrankenPHP}{s7-log}

% ============================================================
%  ALUR DATA
% ============================================================
\section{Alur Data dari Vue ke Laravel ke Basis Data}

\begin{enumerate}[label=\arabic*.]
    \item User membuka browser dan mengakses \texttt{http://localhost:5173}. Request ini dilayani oleh container \texttt{frontend} (Nginx) yang menyajikan file statis hasil build Vue.
    \item User mengisi form login. Vue mengirim \texttt{POST /api/login} ke backend melalui Axios. Request ini diteruskan oleh konfigurasi Nginx ke container \texttt{backend} melalui jaringan internal Docker Compose (menggunakan nama service \texttt{backend}, bukan IP).
    \item Container \texttt{backend} (Laravel + FrankenPHP) menerima request, memvalidasi kredensial terhadap data di PostgreSQL, membuat token Sanctum, dan mengembalikannya dalam response JSON.
    \item Vue menyimpan token di \texttt{localStorage}. Axios interceptor secara otomatis menyertakan token di header \texttt{Authorization: Bearer <token>} pada setiap request berikutnya.
    \item Saat user melakukan operasi CRUD (misal membuat project baru), Axios mengirim \texttt{POST /api/projects} dengan data JSON ke backend.
    \item Laravel memvalidasi input (mengembalikan 422 jika gagal), lalu menyimpan data ke PostgreSQL melalui Eloquent ORM. Data disimpan di named volume \texttt{pgdata} sehingga persisten.
    \item Laravel mengembalikan response JSON (status 201 untuk create, 200 untuk operasi lain).
    \item Vue menerima response dan memperbarui tampilan secara reaktif.
\end{enumerate}

% ============================================================
%  KENDALA
% ============================================================
\section{Kendala dan Cara Mengatasinya}

\begin{longtable}{@{}p{0.42\textwidth}p{0.53\textwidth}@{}}
\toprule
\textbf{Kendala} & \textbf{Solusi} \\
\midrule
\endhead
% Isi dengan kendala nyata selama pengerjaan -- update sesuai pengalaman
CORS error saat frontend Vue mengakses backend Laravel dari container berbeda &
Menambahkan konfigurasi CORS di \texttt{config/cors.php} Laravel dengan \texttt{allowed\_origins} yang mencakup alamat frontend, serta memastikan \texttt{supports\_credentials} diatur \texttt{true} \\[8pt]

PostgreSQL container belum siap saat Laravel mencoba migrasi &
Menambahkan \texttt{depends\_on} dengan \texttt{condition: service\_healthy} dan healthcheck pada service \texttt{db} di docker-compose.yml, sehingga backend menunggu database benar-benar siap \\[8pt]

Path proyek terlalu panjang di Windows sehingga file di \texttt{vendor/} terpotong saat \texttt{composer install} (beberapa file kehilangan ekstensi \texttt{.php}) &
Menggunakan \texttt{subst X: \textless{}path\textgreater{}} untuk memetakan path panjang ke huruf drive pendek, kemudian menjalankan ulang \texttt{composer install} dari drive \texttt{X:\textbackslash{}} \\
\bottomrule
\end{longtable}

% ============================================================
%  KESIMPULAN
% ============================================================
\section{Kesimpulan}

Berdasarkan pengerjaan ujian tengah semester ini, dapat disimpulkan bahwa:

\begin{enumerate}[label=\arabic*.]
    \item Aplikasi \textit{Project Tracker} berhasil dibangun dengan arsitektur backend-frontend terpisah menggunakan Laravel 11 (RESTful API + Sanctum) dan Vue 3 (SPA + Axios + Pinia), saling berkomunikasi melalui API JSON dengan autentikasi token.
    \item Seluruh kode dikelola menggunakan Git dengan strategi branching (\texttt{main}, \texttt{dev}, \texttt{feature/*}), branch protection pada \texttt{main} dan \texttt{dev}, dan Pull Request wajib sebelum merge.
    \item Pipeline CI/CD diterapkan melalui GitHub Actions: job \texttt{test} (backend) dan \texttt{build} (frontend) berjalan pada setiap push/PR, sedangkan job \texttt{deploy} (build dan push image Docker ke GHCR dengan tag commit SHA) hanya berjalan pada push ke \texttt{main}.
    \item Dockerfile backend memanfaatkan FrankenPHP sebagai server aplikasi dengan layer caching Composer, sedangkan Dockerfile frontend menggunakan multi-stage build (Node.js + Nginx) sehingga image akhir tidak mengandung \texttt{node\_modules}.
    \item Docker Compose berhasil mengorkestrasi tiga service (PostgreSQL, backend, frontend) dalam satu perintah \texttt{docker compose up -d --build}, dengan named volume \texttt{pgdata} untuk persistensi data dan komunikasi antar-service melalui jaringan Docker.
    \item Penggunaan FrankenPHP sebagai server aplikasi Laravel (bonus) terbukti berfungsi, dibuktikan dari log container yang menunjukkan request dilayani oleh FrankenPHP.
\end{enumerate}

\end{document}
